\documentclass[11pt]{article}
\usepackage[a4paper,margin=17.5mm]{geometry}
\usepackage[no-math]{fontspec}
\setmainfont{Latin Modern Roman}
\usepackage{amsmath,amssymb,amsthm,xfrac,xspace,hyperref,graphicx,comment}
\excludecomment{DefTralics}
\providecommand{\backmatter}{}

\providecommand{\bibliofont}{\small}
\newcommand{\ie}{i.e.,\xspace}
\newcommand{\eg}{e.g.,\xspace}
\newcommand{\etc}{etc.\xspace}
\newcommand{\cf}{cf.\xspace}
\newcommand{\resp}{resp.\xspace}
\newcommand{\smaller}{\small}
\newcommand{\Subsection}{\subsection}
\newcommand{\Subsubsection}{\subsubsection}
\newcommand{\skpt}{\smallskip}
\emergencystretch=3em
\providecommand{\Mc}{}\input{author-preamble.tex}
\makeatletter
\crefname{corollary}{corollary}{corollaries}
\Crefname{corollary}{Corollary}{Corollaries}
\crefname{definition}{definition}{definitions}
\Crefname{definition}{Definition}{Definitions}
\crefname{example}{example}{examples}
\Crefname{example}{Example}{Examples}
\crefname{lemma}{lemma}{lemmas}
\Crefname{lemma}{Lemma}{Lemmas}
\crefname{proposition}{proposition}{propositions}
\Crefname{proposition}{Proposition}{Propositions}
\crefname{remark}{remark}{remarks}
\Crefname{remark}{Remark}{Remarks}
\crefname{theorem}{theorem}{theorems}
\Crefname{theorem}{Theorem}{Theorems}
\expandafter\def\csname jeplabeltype@thmintro\endcsname{theorem}
\expandafter\def\csname jeplabeltype@rmq:MtimesAHunital\endcsname{remark}
\expandafter\def\csname jeplabeltype@rmq:HChS\endcsname{remark}
\expandafter\def\csname jeplabeltype@thm:LQT\endcsname{theorem}
\expandafter\def\csname jeplabeltype@thm:wodzicki\endcsname{theorem}
\expandafter\def\csname jeplabeltype@lem:filtrationF\endcsname{lemma}
\expandafter\def\csname jeplabeltype@cor:HIHA\endcsname{corollary}
\expandafter\def\csname jeplabeltype@cor:BmodHunital\endcsname{corollary}
\expandafter\def\csname jeplabeltype@lem:kernelsfiltration\endcsname{lemma}
\expandafter\def\csname jeplabeltype@cor:HAHB\endcsname{corollary}
\expandafter\def\csname jeplabeltype@rmk:relKnc\endcsname{remark}
\expandafter\def\csname jeplabeltype@prop:malcev\endcsname{proposition}
\expandafter\def\csname jeplabeltype@thm:goodwillie\endcsname{theorem}
\expandafter\def\csname jeplabeltype@rmq:relChern\endcsname{remark}
\expandafter\def\csname jeplabeltype@cor:goodwillieHunital\endcsname{corollary}
\expandafter\def\csname jeplabeltype@df:schlessinger\endcsname{definition}
\expandafter\def\csname jeplabeltype@ex:artinstacksareFMP\endcsname{example}
\expandafter\def\csname jeplabeltype@thm:pridhamlurie\endcsname{theorem}
\expandafter\def\csname jeplabeltype@ex:tgtcoincides\endcsname{example}
\expandafter\def\csname jeplabeltype@lem:BsimCDB\endcsname{lemma}
\expandafter\def\csname jeplabeltype@df:FMPC\endcsname{definition}
\expandafter\def\csname jeplabeltype@lem:iCleftadjoint\endcsname{lemma}
\expandafter\def\csname jeplabeltype@lem:functoriality\endcsname{lemma}
\expandafter\def\csname jeplabeltype@prop:abelianlie\endcsname{proposition}
\expandafter\def\csname jeplabeltype@lem:abelianization\endcsname{lemma}
\expandafter\def\csname jeplabeltype@prop:calceAb\endcsname{proposition}
\expandafter\def\csname jeplabeltype@prop:fmpQ\endcsname{proposition}
\expandafter\def\csname jeplabeltype@eq:eAb\endcsname{proposition}
\expandafter\def\csname jeplabeltype@rmq:lbar\endcsname{remark}
\expandafter\def\csname jeplabeltype@lem:generateurs\endcsname{lemma}
\expandafter\def\csname jeplabeltype@lem:tgtSQ\endcsname{lemma}
\expandafter\def\csname jeplabeltype@prop:monoidal\endcsname{proposition}
\expandafter\def\csname jeplabeltype@lem:Cactionell\endcsname{lemma}
\expandafter\def\csname jeplabeltype@cor:tensorellzero\endcsname{corollary}
\expandafter\def\csname jeplabeltype@thm:excision\endcsname{theorem}
\expandafter\def\csname jeplabeltype@cor:tgtK\endcsname{corollary}
\expandafter\def\csname jeplabeltype@rmk:ellKnc\endcsname{remark}
\expandafter\def\csname jeplabeltype@rmk:goodwilliederivative\endcsname{remark}
\expandafter\def\csname jeplabeltype@lem:ellBH\endcsname{lemma}
\expandafter\def\csname jeplabeltype@prop:alphaelleq\endcsname{proposition}
\expandafter\def\csname jeplabeltype@lem:gammaeqs\endcsname{lemma}
\expandafter\def\csname jeplabeltype@lem:deltaeqs\endcsname{lemma}
\expandafter\def\csname jeplabeltype@lem:tgtBGL\endcsname{lemma}
\expandafter\def\csname jeplabeltype@thm:compLQT\endcsname{theorem}
\expandafter\def\csname jeplabeltype@rmq:Risk\endcsname{remark}
\AtBeginDocument{\let\jeporiginalLabel\label
\renewcommand{\label}[1]{\ifcsname jeplabeltype@#1\endcsname
\edef\jeplabeltype{\csname jeplabeltype@#1\endcsname}\expandafter\jeporiginalLabel\expandafter[\jeplabeltype]{#1}\else\jeporiginalLabel{#1}\fi}}
\makeatother
\begin{document}
\begin{center}{\Large Stable item 2055}\end{center}
\paragraph{Source and attribution.}
Benjamin Hennion, \emph{The tangent complex of K-theory}, Journal de l’École polytechnique — Mathématiques 8 (2021), publisher version, DOI 10.5802/jep.161. \href{https://jep.centre-mersenne.org/articles/10.5802/jep.161/}{Publisher article}. Authors' text: \href{https://creativecommons.org/licenses/by/4.0/}{CC BY 4.0}.
\paragraph{Editorial problem boundary.}
Prove only Theorem 3.1.1(b) below over the printed characteristic-zero base field, for an algebra object \(\mathcal A\) in the rational pre-formal-moduli category and its pointed fibre \(\overline{\mathcal A}\), assuming the induced nonunital algebra \(A=\ell^{\mathbf Q}(\mathcal A)\) is H-unital. The exact formalised Hochschild and cyclic homology excision equivalences, including the conclusion for the \(\beta\)-maps, are selected. Both needed bar/Hochschild filtrations and nonunital monoidality context are retained. Formal tangent-K-theory corollaries and unproved Goodwillie-derivative remarks remain surrounding context only.
\paragraph{Difficulty (editorial): H3.}
Ordinary Wodzicki excision does not identify these maps after rational formalisation of a varying pre-moduli algebra. The author constructs and compares the relevant bar and Hochschild filtrations, proves that their successive cofibers vanish under the formalisation functor using H-unitality, and supplies the pointed nonunital tensor-compatibility needed in that argument. These steps establish the alpha equivalences and, through the separately complete composite-map argument, the claimed beta equivalences.
\paragraph{Editorial transcription note.}
Original author language, mathematics and ordering are preserved. Publisher front matter is replaced by this independent article wrapper. Bibliography is recovered from matching cached publisher HTML, including its TeX formulas. No new proof or mathematical repair is supplied. Surrounding original results are retained for conservative context and are not extra selected corpus claims.
\clearpage
\input{author-body.tex}
\begin{thebibliography}{99}
\bibitem{beilinson:perv} A. A. Beilinson - “On the derived category of perverse sheaves” , in $K$-theory, arithmetic and geometry (Moscow, 1984–1986) , Lect. Notes in Math. , vol. 1289 , Springer, Berlin, 1987, p. 27-41
\bibitem{beilinson:relK} A. A. Beilinson - “Relative continuous $K$-theory and cyclic homology” , Münster J. Math. 7 (2014) no. 1, p. 51-81
\bibitem{blanckatzarkovpandit} A. Blanc, L. Katzarkov \& P. Pandit - “Generators in formal deformations of categories” , Compositio Math. 154 (2018) no. 10, p. 2055-2089
\bibitem{bloch:tangentk} S. Bloch - “On the tangent space to Quillen $K$-theory” , in Algebraic $K$-theory, I: Higher $K$-theories (Proc. Conf., Battelle Memorial Inst., Seattle, Wash., 1972) , Lect. Notes in Math. , vol. 341 , Springer, 1973, p. 205-210
\bibitem{burghelea:cyclicK} D. Burghelea - “Cyclic homology and the algebraic $K$-theory of spaces. I” , in Applications of algebraic $K$-theory to algebraic geometry and number theory, Part I, II (Boulder, Colo., 1983) , Contemp. Math. , vol. 55 , American Mathematical Society, Providence, RI, 1986, p. 89-115
\bibitem{cathelineau:lambda} J.-L. Cathelineau - “$\lambda $-structures in algebraic $K$-theory and cyclic homology” , $K$-Theory 4 (1990/91) no. 6, p. 591-606
\bibitem{chw:lambda} G. Cortiñas, C. Haesemeyer \& C. A. Weibel - “Infinitesimal cohomology and the Chern character to negative cyclic homology” , Math. Ann. 344 (2009) no. 4, p. 891-922
\bibitem{cortinasweibel:relative} G. Cortiñas \& C. A. Weibel - “Relative Chern characters for nilpotent ideals” , in Algebraic topology , Abel Symp. , vol. 4 , Springer, Berlin, 2009, p. 61-82
\bibitem{DonadzeLadra} G. Donadze \& M. Ladra - “The excision theorems in Hochschild and cyclic homologies” , Proc. Roy. Soc. Edinburgh Sect. A 144 (2014) no. 2, p. 305-317
\bibitem{faontehennionkapranov:kacmoody} G. Faonte, B. Hennion \& M. Kapranov - “Higher Kac-Moody algebras and moduli spaces of $G$-bundles” , Adv. Math. 346 (2019), p. 389-466
\bibitem{feigintsygan:addK} B. L. Feigin \& B. L. Tsygan - “Additive $K$-theory” , in $K$-theory, arithmetic and geometry (Moscow, 1984–1986) , Lect. Notes in Math. , vol. 1289 , Springer, Berlin, 1987, p. 67-209
\bibitem{guccioneguccione} J. A. Guccione \& J. J. Guccione - “The theorem of excision for Hochschild and cyclic homology” , J. Pure Appl. Algebra 106 (1996) no. 1, p. 57-60
\bibitem{goodwillie:tangentk} T. G. Goodwillie - “Relative algebraic $K$-theory and cyclic homology” , Ann. of Math. (2) 124 (1986) no. 2, p. 347-402
\bibitem{loday:hc} J.-L. Loday - Cyclic homology , Grundlehren Math. Wiss. , vol. 301 , Springer-Verlag, Berlin, 1992
\bibitem{lodayquillen} J.-L. Loday \& D. Quillen - “Homologie cyclique et homologie de l’algèbre de Lie des matrices” , C. R. Acad. Sci. Paris Sér. I Math. 296 (1983) no. 6, p. 295-297
\bibitem{lurie:htt} J. Lurie - Higher topos theory , Annals of Math. Studies , vol. 170 , Princeton University Press, Princeton, NJ, 2009
\bibitem{lurie:dagx} J. Lurie - “Derived algebraic geometry X: Formal moduli problems” (2011), available at https://www.math.ias.edu/\textasciitilde{}lurie/papers/DAG-X.pdf
\bibitem{lurie:halg} J. Lurie - “Higher algebra” (2016), available at https://www.math.ias.edu/\textasciitilde{}lurie/papers/HA.pdf
\bibitem{malcev:lie} A. I. Mal’cev - “Nilpotent torsion-free groups” , Izv. Akad. Nauk SSSR Ser. Mat. 13 (1949), p. 201-212
\bibitem{pridham:deformation} J. P. Pridham - “Unifying derived deformation theories” , Adv. Math. 224 (2010) no. 3, p. 772-826, Corrigendum: Ibid. 228 (2011), no. 4, p. 2554–2556
\bibitem{pridham:smoothfunctions} J. P. Pridham - “Smooth functions on algebraic K-theory” , 2016
\bibitem{schwedeshipley:monoidalDK} S. Schwede \& B. Shipley - “Equivalences of monoidal model categories” , Algebraic Geom. Topol. 3 (2003), p. 287-334
\bibitem{suslin:eqK} A. A. Suslin - “On the equivalence of $K$-theories” , Comm. Algebra 9 (1981) no. 15, p. 1559-1566
\bibitem{suslinwodzicki} A. A. Suslin \& M. Wodzicki - “Excision in algebraic $K$-theory” , Ann. of Math. (2) 136 (1992) no. 1, p. 51-122
\bibitem{tamme:excisionrevisited} G. Tamme - “Excision in algebraic $K$-theory revisited” , Compositio Math. 154 (2018) no. 9, p. 1801-1814
\bibitem{tsygan:HC} B. L. Tsygan - “Homology of matrix Lie algebras over rings and the Hochschild homology” , Uspehi Mat. Nauk 38 (1983) no. 2(230), p. 217-218
\bibitem{toen:hagii} B. Toën \& G. Vezzosi - Homotopical algebraic geometry. II. Geometric stacks and applications , Mem. Amer. Math. Soc. , vol. 193, no. 902 , American Mathematical Society, Providence, RI, 2008
\bibitem{volodin:ktheory} I. A. Volodin - “Algebraic $K$-theory as an extraordinary homology theory on the category of associative rings with a unit” , Izv. Akad. Nauk SSSR Ser. Mat. 35 (1971), p. 844-873
\bibitem{waldhausen:ktheory} F. Waldhausen - “Algebraic $K$-theory of spaces” , in Algebraic and geometric topology (New Brunswick, N.J., 1983) , Lect. Notes in Math. , vol. 1126 , Springer, Berlin, 1985, p. 318-419
\bibitem{weibel:cyclichodge} C. A. Weibel - “The Hodge filtration and cyclic homology” , $K$-Theory 12 (1997) no. 2, p. 145-164
\bibitem{wodzicki:excision} M. Wodzicki - “Excision in cyclic homology and in rational algebraic $K$-theory” , Ann. of Math. (2) 129 (1989) no. 3, p. 591-639
\end{thebibliography}
\end{document}
